% ============================================================================
\section{Tau be continued}
\label{sec:future}
% ============================================================================
Four directions follow from the remaining accuracy--resource trade-off:
\begin{itemize}
  \item \textbf{Smaller CNNs.} Reduce channels and head widths, prune, quantise and distil strong models into compact students; validate the resulting efficiency, latency and FPGA resource use.
  \item \textbf{Adaptive loss design.} Derive target regions and tolerances from the starting model's weaknesses, with uncertainty-aware updates that adapt to the input representation.
  \item \textbf{Whole-event learning.} Vary the number of input objects and design an event-level architecture that can learn dependencies between candidates, rather than scoring each in isolation.
  \item \textbf{Richer detector inputs.} The limited detector information available to our models leaves low-energy tau and background deposits difficult to distinguish. Richer measurements and finer spatial resolution at the input stage could therefore be crucial, rather than further model complexity alone.
\end{itemize}
\Cref{app:future} develops these directions and their evaluation requirements.
